\chapter{Connecting external agents through MCP}\label{ch:mcp}
\chapterintro{The tool contract is exposed as a Model Context Protocol server, \code{agentto.mcp\_server}. Any MCP
client can drive a design session with it: attached to the live workspace, where its edits appear as you watch, or in a
session of its own. This chapter covers the server's modes and profiles, configuration for Claude Code and the Codex
CLI, and the conventions an agent needs to know.}

\section{The server}
\index{MCP server}
The server speaks MCP over standard input and output and publishes one MCP tool per contract tool, with the tool's
docstring as its description and its signature as the argument schema. Every call is forwarded to a
\code{Session} and validated there, so a malformed call is rejected and logged like any other failed call, and the agent
receives the error text. Results are compact JSON text; \tool{render} returns its picture as an image.

The server runs in one of two modes:
\begin{description}
  \item[Attach] (\code{--attach http://127.0.0.1:8765}). The agent drives the session the workspace has open. Its calls go
  to the backend's \code{/api/tools/}\meta{name} endpoint with actor \code{agent}, appear in the workspace as they happen,
  and can be undone like any other step. Solves run on the workspace's GPU worker.
  \item[Stand-alone] (\code{--session DIR}). The server keeps its own session in a directory, with its own log and its
  own GPU worker. This is how the paper's benchmarks ran.
\end{description}

\subsection*{Command-line flags}
% Generated by docs/guidebook/gen_tables.py from the code. Do not edit by hand.
\begin{reftable}{@{}>{\raggedright\arraybackslash}p{0.2\linewidth}>{\raggedright\arraybackslash}p{0.18\linewidth}X@{}}
Flag & Values & Meaning \\ \midrule
\code{-{}-session}\index{session@\texttt{-{}-session} (MCP flag)} & directory & session directory (stand-alone mode) \\
\code{-{}-attach}\index{attach@\texttt{-{}-attach} (MCP flag)} & URL & backend URL of a running GUI session \\
\code{-{}-profile}\index{profile@\texttt{-{}-profile} (MCP flag)} & \code{spec} | \code{clarify} | \code{spec\_\allowbreak{}prov} | \code{clarify\_\allowbreak{}prov} | \code{full} (default \code{spec}) & which tool set the server offers (the profile table above) \\
\code{-{}-seed}\index{seed@\texttt{-{}-seed} (MCP flag)} & JSON file & stand-alone mode: JSON spec the session starts from (logged as actor system) \\
\code{-{}-oracle}\index{oracle@\texttt{-{}-oracle} (MCP flag)} & JSON file & stand-alone mode: intended-reading JSON; a scripted oracle answers ask\_clarification (X2 condition G+C) \\
\code{-{}-answer-url}\index{answer-url@\texttt{-{}-answer-url} (MCP flag)} & URL & stand-alone mode: a GUI backend that puts ask\_clarification to its user \\
\code{-{}-no-rules}\index{no-rules@\texttt{-{}-no-rules} (MCP flag)} & JSON file & stand-alone mode: validate\_spec checks the schema only, no rule findings (X2 ablation GNR) \\
\end{reftable}


\section{Profiles}\label{sec:profiles}
\index{profile}
A profile decides which tools the server offers. The default is \code{spec}; give \code{--profile full} for an agent
that should solve and export.

% Generated by docs/guidebook/gen_tables.py from the code. Do not edit by hand.
\begin{reftable}{@{}lc>{\raggedright\arraybackslash}X@{}}
Profile & Tools & Contents \\ \midrule
\code{spec} & 16 & \code{new\_\allowbreak{}spec}, \code{set\_\allowbreak{}spec}, \code{set\_\allowbreak{}params}, \code{add\_\allowbreak{}support}, \code{add\_\allowbreak{}load}, \code{add\_\allowbreak{}passive}, \code{update\_\allowbreak{}item}, \code{remove\_\allowbreak{}item}, \code{undo}, \code{redo}, \code{get\_\allowbreak{}spec}, \code{list\_\allowbreak{}regions}, \code{validate\_\allowbreak{}spec}, \code{explain\_\allowbreak{}spec}, \code{render}, \code{ladder} \\
\code{clarify} & 17 & everything in \code{spec}, plus \code{ask\_\allowbreak{}clarification} \\
\code{full} & 27 & everything in \code{clarify}, plus \code{load\_\allowbreak{}template}, \code{set\_\allowbreak{}source}, \code{list\_\allowbreak{}sources}, \code{solve}, \code{status}, \code{wait\_\allowbreak{}job}, \code{cancel}, \code{field\_\allowbreak{}stats}, \code{export\_\allowbreak{}stl}, \code{export\_\allowbreak{}report} \\
\end{reftable}


In the paper, \code{spec} was the grounding condition G, \code{clarify} the condition G+C and the Propose assistant, and
\code{full} the complete sessions from brief to release, Act mode, and attaching to the workspace. The follow-up
conditions used two further profiles. \code{spec\_prov} adds \tool{set_source} and \tool{list_sources} to \code{spec}
(condition G+P): the agent tags each specification item as stated, derived, assumed or answered.
\code{clarify\_prov} does the same for \code{clarify} and limits questions to items tagged assumed (condition G+C+P).
Under both, \tool{validate_spec} adds an advisory note with code R014 that lists the items tagged assumed and the
untagged items; it never blocks a solve and is not one of the thirteen rules. The \code{-{}-no-rules} flag gave the
rule-feedback ablation GNR.

\section{Attach Claude Code to the workspace}
With the backend running, save this as \code{mcp.json}, replacing the two paths with your Python interpreter of the
\code{agentto} environment and the repository's \code{src} directory:

\begin{jsoncode}
{"mcpServers": {"agentto": {
  "command": "/path/to/envs/agentto/bin/python",
  "args": ["-B", "-m", "agentto.mcp_server", "--profile", "full",
           "--attach", "http://127.0.0.1:8765"],
  "env": {"PYTHONPATH": "/path/to/AgenticTO/src"}}}}
\end{jsoncode}

\begin{shell}
claude --mcp-config mcp.json
\end{shell}

Allow the \code{mcp\_\_agentto} tools when Claude Code asks, or start it with \code{--allowedTools mcp\_\_agentto}.
Set \code{AGENTTO\_VOXFEM\_SRC} in the environment block too if the server runs stand-alone and must solve.

In Figure~\ref{fig:attach}, Claude Code (\code{claude-opus-5-5} on a claude.ai subscription) was asked to add a 300~N
load on the top face of the part from Chapter~\ref{ch:assistant} and to validate the specification. It took 13~s. Its
calls appear in the timeline as \gui{AGENT} records, and the new load appears in the tree and the viewport.

\begin{figure}[t]
\wideshot{26_attach.png}
\caption{Claude Code, attached over MCP, added \code{load\_claude} to the live design.}
\label{fig:attach}
\end{figure}

\begin{figure}[h]
\widegraphic[178mm]{27_attach_timeline.png}
\caption{The timeline after the Claude Code session: its tool calls are tagged \gui{AGENT}.}
\label{fig:attachtimeline}
\end{figure}

\needspace{14\baselineskip}
\section{Attach the Codex CLI}
Add this to \code{\textasciitilde/.codex/config.toml}:

\begin{tomlcode}
[mcp_servers.agentto]
command = "/path/to/envs/agentto/bin/python"
args = ["-B", "-m", "agentto.mcp_server", "--profile", "full",
        "--attach", "http://127.0.0.1:8765"]
env = { PYTHONPATH = "/path/to/AgenticTO/src" }
startup_timeout_sec = 60
tool_timeout_sec = 900
default_tools_approval_mode = "approve"
\end{tomlcode}

\begin{warnbox}[Approval in non-interactive runs]
Without \code{default\_tools\_approval\_mode = "approve"}, non-interactive Codex runs (\code{codex exec}) reject every MCP
call, because they cannot ask you. The AgenticTO tools only touch the session directory. A \code{tool\_timeout\_sec}
shorter than your longest solve cuts the call off; the solve itself keeps running, and the agent can collect it with
\tool{wait_job}.
\end{warnbox}

\needspace{9\baselineskip}
\section{Stand-alone sessions}
Without the workspace, the server keeps its own session:

\begin{shell}
python -B -m agentto.mcp_server --session runs/my_session --profile full
python -B -m agentto.mcp_server --session runs/b01 --profile clarify --seed spec.json
python -B -m agentto.mcp_server --session runs/b01 --profile clarify \
    --answer-url http://127.0.0.1:8765
\end{shell}

\code{--seed} starts the session from an existing specification (logged as actor \code{system}). \code{--answer-url}
sends the agent's clarification questions to the user of a running workspace; this is how the Propose assistant asks
you its questions. \code{--oracle} makes a scripted oracle answer them from a reference reading (the benchmark's
condition G+C); without an answerer, a question returns at once with the advice to choose the most reasonable reading
and state it as an assumption.

\section{Conventions an agent needs}\label{sec:agentconventions}
The server's instructions and the tool descriptions carry most of these; they are listed here so that you can write
your own prompts.

\begin{description}
  \item[Coordinates.] Absolute, origin at the domain's minimum corner. Patch centres and sizes are in the face's two
  in-plane axes in increasing order. \tool{new_spec} and \tool{list_regions} return the in-plane axes of every face.
  \item[Grounding.] After adding an item, \tool{list_regions} reports how many nodes or elements it selects; an item that
  selects 0 does nothing. \tool{validate_spec} returns the rule findings. \tool{render} returns a picture to check
  placement and directions.
  \item[Long solves.] \tool{solve} blocks until the job finishes or \code{timeout\_s} (default \SolveTimeout~s) passes.
  On a timeout it returns the job id with a hint; the agent should call \tool{wait_job}, not cancel and resubmit, which
  would start the run again from zero. \code{wait=False} returns the job id at once; \tool{status} lists the jobs.
  \item[Rungs.] Ask for a rung by name. The release rung is often finer than the specification's \code{h}, which still
  satisfies a brief that asks for elements of at most \code{h}. Asking for a rung and an \code{h} that disagree is an
  error.
  \item[Certificates.] The agent can read the certificate in every solve result, but no tool can write one. Only a
  certified release-rung result is a released design; anything else is exported as a draft.
  \item[Clarification.] One question at a time, with two to four options, each a complete item or parameter values.
  The call returns once the user has answered.
\end{description}

The complete tool reference, with every argument, its type and its default, is Appendix~\ref{app:tools}. The summary
below shows which tools change the specification, which profiles offer them, and how replay treats them.

% Generated by docs/guidebook/gen_tables.py from the code. Do not edit by hand.
\begin{reftable}{@{}>{\raggedright\arraybackslash}p{0.22\linewidth}lccc>{\raggedright\arraybackslash}X@{}}
Tool & Changes spec & \code{spec} & \code{clarify} & \code{full} & Replay \\ \midrule
\multicolumn{6}{@{}l}{\textsf{\bfseries Build the spec}} \\
\hyperref[tool:new_spec]{\code{new\_\allowbreak{}spec}} & yes & \yes & \yes & \yes & re-run, compared \\
\hyperref[tool:set_spec]{\code{set\_\allowbreak{}spec}} & yes & \yes & \yes & \yes & re-run, compared \\
\hyperref[tool:load_template]{\code{load\_\allowbreak{}template}} & yes & \no & \no & \yes & re-run, compared \\
\hyperref[tool:set_params]{\code{set\_\allowbreak{}params}} & yes & \yes & \yes & \yes & re-run, compared \\
\hyperref[tool:add_support]{\code{add\_\allowbreak{}support}} & yes & \yes & \yes & \yes & re-run, compared \\
\hyperref[tool:add_load]{\code{add\_\allowbreak{}load}} & yes & \yes & \yes & \yes & re-run, compared \\
\hyperref[tool:add_passive]{\code{add\_\allowbreak{}passive}} & yes & \yes & \yes & \yes & re-run, compared \\
\hyperref[tool:update_item]{\code{update\_\allowbreak{}item}} & yes & \yes & \yes & \yes & re-run, compared \\
\hyperref[tool:remove_item]{\code{remove\_\allowbreak{}item}} & yes & \yes & \yes & \yes & re-run, compared \\
\hyperref[tool:undo]{\code{undo}} & yes & \yes & \yes & \yes & re-run, compared \\
\hyperref[tool:redo]{\code{redo}} & yes & \yes & \yes & \yes & re-run, compared \\
\multicolumn{6}{@{}l}{\textsf{\bfseries Inspect and check}} \\
\hyperref[tool:get_spec]{\code{get\_\allowbreak{}spec}} & -- & \yes & \yes & \yes & re-run, compared \\
\hyperref[tool:list_regions]{\code{list\_\allowbreak{}regions}} & -- & \yes & \yes & \yes & re-run, compared \\
\hyperref[tool:validate_spec]{\code{validate\_\allowbreak{}spec}} & -- & \yes & \yes & \yes & re-run, compared \\
\hyperref[tool:explain_spec]{\code{explain\_\allowbreak{}spec}} & -- & \yes & \yes & \yes & re-run, compared \\
\hyperref[tool:render]{\code{render}} & -- & \yes & \yes & \yes & re-run, compared \\
\hyperref[tool:ladder]{\code{ladder}} & -- & \yes & \yes & \yes & re-run, compared \\
\multicolumn{6}{@{}l}{\textsf{\bfseries Ask the user}} \\
\hyperref[tool:ask_clarification]{\code{ask\_\allowbreak{}clarification}} & -- & \no & \yes & \yes & re-run, compared \\
\multicolumn{6}{@{}l}{\textsf{\bfseries Record value sources}} \\
\hyperref[tool:set_source]{\code{set\_\allowbreak{}source}} & -- & \no & \no & \yes & re-run, compared \\
\hyperref[tool:list_sources]{\code{list\_\allowbreak{}sources}} & -- & \no & \no & \yes & re-run, compared \\
\multicolumn{6}{@{}l}{\textsf{\bfseries Solve}} \\
\hyperref[tool:solve]{\code{solve}} & -- & \no & \no & \yes & re-solved, compared \\
\hyperref[tool:status]{\code{status}} & -- & \no & \no & \yes & skipped \\
\hyperref[tool:wait_job]{\code{wait\_\allowbreak{}job}} & -- & \no & \no & \yes & skipped \\
\hyperref[tool:cancel]{\code{cancel}} & -- & \no & \no & \yes & skipped \\
\hyperref[tool:field_stats]{\code{field\_\allowbreak{}stats}} & -- & \no & \no & \yes & re-run, compared \\
\multicolumn{6}{@{}l}{\textsf{\bfseries Export}} \\
\hyperref[tool:export_stl]{\code{export\_\allowbreak{}stl}} & -- & \no & \no & \yes & re-run, compared \\
\hyperref[tool:export_report]{\code{export\_\allowbreak{}report}} & -- & \no & \no & \yes & re-run, compared \\
\end{reftable}


\section{The HTTP interface}\label{sec:http}
\index{HTTP API}
The backend's own endpoints are listed below, generated from \code{src/agentto/app/server.py}. Attach mode uses only
\code{POST /api/tools/}\meta{name} with the body \code{\{"args": \{\ldots\}, "actor": "agent"\}}; the workspace uses the
WebSocket. The backend listens on \code{127.0.0.1} only.

% Generated by docs/guidebook/gen_tables.py from the code. Do not edit by hand.
\begin{reftable}{@{}l>{\raggedright\arraybackslash}p{0.3\linewidth}X@{}}
Method & Path & Purpose (from the handler's docstring or name) \\ \midrule
GET & \code{/\allowbreak{}api/\allowbreak{}health} & Version and schema version. \\
GET & \code{/\allowbreak{}api/\allowbreak{}schema} & The live JSON Schema of the spec. \\
GET & \code{/\allowbreak{}api/\allowbreak{}library} & Names of the template problems. \\
GET & \code{/\allowbreak{}api/\allowbreak{}library/\allowbreak{}\{name\}} & A template problem as a spec (query parameter k: cells on the short edge). \\
GET & \code{/\allowbreak{}api/\allowbreak{}tasks} & Draft X2 tasks (briefs/intended/*.json), for development: their intended specs load like library items. \\
GET & \code{/\allowbreak{}api/\allowbreak{}tasks/\allowbreak{}\{slug\}} & A suite task: its intended spec, expected response and brief text. \\
POST & \code{/\allowbreak{}api/\allowbreak{}validate} & Run the schema and the rules on a posted spec. \\
POST & \code{/\allowbreak{}api/\allowbreak{}ladder} & The rungs for a posted spec. \\
POST & \code{/\allowbreak{}api/\allowbreak{}tools/\allowbreak{}\{name\}} & Tool contract over HTTP (used by the MCP server in -{}-attach mode). \\
GET & \code{/\allowbreak{}api/\allowbreak{}session} & The live spec, the tail of the session log and the chat. \\
POST & \code{/\allowbreak{}api/\allowbreak{}clarify} & ask\_clarification from a stand-alone (propose-mode) agent session: answered by this app's user. \\
GET & \code{/\allowbreak{}api/\allowbreak{}mesh/\allowbreak{}\{job\}} & Preview surface of a finished design (pooled to at most 1.5M triangles): uint32 n\_vertices, uint32 n\_triangles, float32 h\_eff, then float32 positions (spec units) and uint32 indices. \\
GET & \code{/\allowbreak{}api/\allowbreak{}files/\allowbreak{}\{name\}} & Download an exported file (STL, report) of the live session. \\
WS & \code{/\allowbreak{}ws/\allowbreak{}session} & JSON messages \{type, ...\}: ping, validate, tool \{name, args, meta, ref\}, worker\_start. \\
\end{reftable}

